\documentclass[10pt,a4paper]{amsart}
\usepackage{amsmath,amssymb,amsthm,graphicx}
% Retained minimal preamble for the Stacks Project chapter more-algebra.tex (More on Algebra).
% The chapter file pinned-original-source.tex carries only \input{preamble} on its first line, so no
% preamble text lives inside the chapter itself.  The Stacks Project uses one single shared preamble for
% every chapter; the snapshot stages it identically as sources/stacks-chow/preamble.tex and
% sources/stacks-etale-cohomology/preamble.tex.  The block below is copied mechanically, line for line,
% from sources/stacks-chow/preamble.tex lines 14-18 (the Xy-pic setup) and 154-217 (the native theorem
% declarations together with the mathematical macros).  Nothing was added, removed or redefined; the
% wrapper class and the amsmath/amssymb/amsthm/graphicx base packages are supplied by the assembler.
% Chapter: more-algebra.tex -- More on Algebra.
% For commutative diagrams we use Xy-pic
\usepackage[all]{xy}

% We use 2cell for 2-commutative diagrams.
\xyoption{2cell}

% Theorem environments.
%
\theoremstyle{plain}
\newtheorem{theorem}[subsection]{Theorem}
\newtheorem{proposition}[subsection]{Proposition}
\newtheorem{lemma}[subsection]{Lemma}

\theoremstyle{definition}
\newtheorem{definition}[subsection]{Definition}
\newtheorem{example}[subsection]{Example}
\newtheorem{exercise}[subsection]{Exercise}
\newtheorem{situation}[subsection]{Situation}

\theoremstyle{remark}
\newtheorem{remark}[subsection]{Remark}
\newtheorem{remarks}[subsection]{Remarks}

\numberwithin{equation}{subsection}

% Macros
%
\def\lim{\mathop{\mathrm{lim}}\nolimits}
\def\colim{\mathop{\mathrm{colim}}\nolimits}
\def\Spec{\mathop{\mathrm{Spec}}}
\def\Hom{\mathop{\mathrm{Hom}}\nolimits}
\def\Ext{\mathop{\mathrm{Ext}}\nolimits}
\def\SheafHom{\mathop{\mathcal{H}\!\mathit{om}}\nolimits}
\def\SheafExt{\mathop{\mathcal{E}\!\mathit{xt}}\nolimits}
\def\Sch{\mathit{Sch}}
\def\Mor{\mathop{\mathrm{Mor}}\nolimits}
\def\Ob{\mathop{\mathrm{Ob}}\nolimits}
\def\Sh{\mathop{\mathit{Sh}}\nolimits}
\def\NL{\mathop{N\!L}\nolimits}
\def\CH{\mathop{\mathrm{CH}}\nolimits}
\def\proetale{{pro\text{-}\acute{e}tale}}
\def\etale{{\acute{e}tale}}
\def\QCoh{\mathit{QCoh}}
\def\Ker{\mathop{\mathrm{Ker}}}
\def\Im{\mathop{\mathrm{Im}}}
\def\Coker{\mathop{\mathrm{Coker}}}
\def\Coim{\mathop{\mathrm{Coim}}}

% Boxtimes
%
\DeclareMathSymbol{\boxtimes}{\mathbin}{AMSa}{"02}

%
% Macros for moduli stacks/spaces
%
\def\QCohstack{\mathcal{QC}\!\mathit{oh}}
\def\Cohstack{\mathcal{C}\!\mathit{oh}}
\def\Spacesstack{\mathcal{S}\!\mathit{paces}}
\def\Quotfunctor{\mathrm{Quot}}
\def\Hilbfunctor{\mathrm{Hilb}}
\def\Curvesstack{\mathcal{C}\!\mathit{urves}}
\def\Polarizedstack{\mathcal{P}\!\mathit{olarized}}
\def\Complexesstack{\mathcal{C}\!\mathit{omplexes}}
% \Pic is the operator that assigns to X its picard group, usage \Pic(X)
% \Picardstack_{X/B} denotes the Picard stack of X over B
% \Picardfunctor_{X/B} denotes the Picard functor of X over B
\def\Pic{\mathop{\mathrm{Pic}}\nolimits}
\def\Picardstack{\mathcal{P}\!\mathit{ic}}
\def\Picardfunctor{\mathrm{Pic}}
\def\Deformationcategory{\mathcal{D}\!\mathit{ef}}

% --- appended 2026-10-12 (Chief): the project own definitions of the reference/slogan/history
% display environments, copied mechanically from the shared Stacks preamble lines 8-12
% (\usepackage{verbatim} plus the three \newenvironment lines, which the chapter preamble keeps
% verbatim from the project preamble). Nothing added, removed or redefined: these environments
% carry no mathematical content and the project itself typesets them as comments. Without these
% three lines a counted statement whose own text uses them fails with Environment slogan undefined.
\usepackage{verbatim}
\newenvironment{reference}{\comment}{\endcomment}
\newenvironment{slogan}{\comment}{\endcomment}
\newenvironment{history}{\comment}{\endcomment}
\usepackage{cleveref}
\crefname{originallemma}{original lemma}{original lemmas}
\Crefname{originallemma}{Original lemma}{Original lemmas}
\makeatletter
\newlabel{lemma-lift-acyclic-complex}{{0BC9}{source}{}{sourcepointer-1}{}}
\newlabel{lemma-lift-acyclic-complex@cref}{{[originallemma][0][]0BC9}{[1][0][]source}{}{}{}}
\makeatother
\crefname{originallemma}{original lemma}{original lemmas}
\Crefname{originallemma}{Original lemma}{Original lemmas}
\makeatletter
\newlabel{derived-lemma-morphisms-from-projective-complex}{{064B}{source}{}{sourcepointer-2}{}}
\newlabel{derived-lemma-morphisms-from-projective-complex@cref}{{[originallemma][0][]064B}{[1][0][]source}{}{}{}}
\makeatother
\crefname{originaldefinition}{original definition}{original definitions}
\Crefname{originaldefinition}{Original definition}{Original definitions}
\makeatletter
\newlabel{derived-definition-cone}{{014E}{source}{}{sourcepointer-3}{}}
\newlabel{derived-definition-cone@cref}{{[originaldefinition][0][]014E}{[1][0][]source}{}{}{}}
\makeatother
\crefname{originallemma}{original lemma}{original lemmas}
\Crefname{originallemma}{Original lemma}{Original lemmas}
\makeatletter
\newlabel{derived-lemma-morphisms-lift-projective}{{0649}{source}{}{sourcepointer-4}{}}
\newlabel{derived-lemma-morphisms-lift-projective@cref}{{[originallemma][0][]0649}{[1][0][]source}{}{}{}}
\makeatother
\title{A bounded above complex over the quotient can be lifted to a bounded above complex in a class of projective modules}
\author{The Stacks Project Authors}
\date{stacks-project:more-algebra.tex:local-snapshot-20261012; GNU FDL 1.2 or later}
\begin{document}
\maketitle
\setcounter{subsection}{1}
\noindent\textit{Editorial note.} All delivered mathematical text is the author's own, copied from the pinned Stacks Project chapter file more-algebra.tex; no statement or proof step was rewritten, reordered or removed. The counted span is the author environment and its own proof environment. References inside the span are handed to the assembler as navigation stubs that carry the canonical tag of the cited tag only; whether each of them is load-bearing is recorded on the stub itself.\par\medskip
\noindent\textit{Allowed standard prerequisites (editorial).}\begin{itemize}
\item A bounded above complex of projective modules over the quotient admits a quasi-isomorphism to any representing complex of the derived tensor product. used at proof lines 19411-19417. Produces the map delta from the given complex over R/I to the complex K/IK and thereby the cone on delta.
\item An acyclic bounded above complex whose terms are of the form P/IP for P in the class can be lifted to an acyclic complex with terms in the class. invoked at proof lines 19422-19425. Supplies the acyclic lift A of the cone whose terms lie in the class.
\item A map into a bounded above complex of projectives can be lifted along a quasi-isomorphism, which produces the dotted arrow in the diagram. invoked at proof lines 19431-19434. Gives the map from K to A making the square commute and hence the comparison of terms.
\item A finite free or projective module over the class is projective, so kernels and cokernels appearing in the split argument are again projective. used at proof lines 19445-19460. Lets the maps g and the projection be split, producing the complementary summand T.
\end{itemize}\medskip
\ifdefined\hypertarget\hypertarget{sourcepointer-1}{}\fi
\noindent\textit{Source reference.} Original lemma 0BC9: sources/stacks-more-algebra/pinned-original-source.tex, source lines 19343--19343. Invoked at proof lines 19422-19424 to lift the acyclic cone C to an acyclic complex A with terms in the class; turning the cone into an object the statement applies to, namely writing its terms as P/IP, is done inside the counted span.\par\medskip
\ifdefined\hypertarget\hypertarget{sourcepointer-2}{}\fi
\noindent\textit{Source reference.} Original lemma 064B: sources/stacks-chow/tags.txt, source lines 1--1. Used at proof lines 19414-19416 for the existence of the quasi-isomorphism delta from E to K/IK; the construction of the cone and everything after it is inside the counted span.\par\medskip
\ifdefined\hypertarget\hypertarget{sourcepointer-3}{}\fi
\noindent\textit{Source reference.} Original definition 014E: sources/stacks-chow/tags.txt, source lines 1--1. The identity at proof lines 19418-19422 that the i-th term of the cone is the direct sum of K\textasciicircum{}i/IK\textasciicircum{}i with E\textasciicircum{}\{i+1\} is the only use; the closure argument that puts that term in the class is inside the counted span.\par\medskip
\ifdefined\hypertarget\hypertarget{sourcepointer-4}{}\fi
\noindent\textit{Source reference.} Original lemma 0649: sources/stacks-chow/tags.txt, source lines 1--1. Supplies the dotted arrow at proof lines 19431-19434; the subsequent argument that its components are split injections and the identification of the cokernel as the representing complex are inside the counted span.\par\medskip
\begin{lemma}
\label{lemma-lift-complex}
Let $R$ be a ring. Let $I \subset R$ be an ideal. Let $\mathcal{P}$
be a class of $R$-modules.
Let $K \in D(R)$ and let $E^\bullet$ be a complex of $R/I$-modules
representing $K \otimes_R^\mathbf{L} R/I$. Assume
\begin{enumerate}
\item each $P \in \mathcal{P}$ is a projective $R$-module,
\item $P_1 \in \mathcal{P}$ and $P_1 \oplus P_2 \in \mathcal{P}$
if and only if $P_1, P_2 \in \mathcal{P}$,
\item if $f : P_1 \to P_2$, $P_1, P_2 \in \mathcal{P}$ is surjective
modulo $I$, then $f$ is surjective,
\item $E^\bullet$ is bounded above and $E^i$ is of the form $P/IP$
for $P \in \mathcal{P}$, and
\item $K$ can be represented by a bounded above complex
whose terms are in $\mathcal{P}$.
\end{enumerate}
Then there exists a bounded above complex $P^\bullet$ whose terms
are in $\mathcal{P}$ with $P^\bullet/IP^\bullet$ isomorphic to
$E^\bullet$ and representing $K$ in $D(R)$.
\end{lemma}

\begin{proof}
By assumption (5) we can represent $K$ by a bounded above complex 
$K^\bullet$ whose terms are in $\mathcal{P}$. Then $K \otimes_R^\mathbf{L} R/I$
is represented by $K^\bullet/IK^\bullet$. Since $E^\bullet$ is
a bounded above complex of projective $R/I$-modules by (4), 
we can choose a quasi-isomorphism $\delta : E^\bullet \to K^\bullet/IK^\bullet$
(Derived Categories, Lemma
\ref{derived-lemma-morphisms-from-projective-complex}).
Let $C^\bullet$ be cone on $\delta$
(Derived Categories, Definition \ref{derived-definition-cone}).
The module $C^i$ is the direct sum $K^i/IK^i \oplus E^{i + 1}$
hence is of the form $P/IP$ for some $P \in \mathcal{P}$
as (2) says in particular that $\mathcal{P}$ is preserved under taking sums.
Since $C^\bullet$ is acyclic, we can apply
Lemma \ref{lemma-lift-acyclic-complex} and find a acyclic
lift $A^\bullet$ of $C^\bullet$. The complex $A^\bullet$ is
bounded above and has terms in $\mathcal{P}$. In
$$
\xymatrix{
K^\bullet \ar@{..>}[r] \ar[d] & A^\bullet \ar[d] \\
K^\bullet/IK^\bullet \ar[r] & C^\bullet \ar[r] & E^\bullet[1]
}
$$
we can find the dotted arrow making the diagram commute
by Derived Categories, Lemma
\ref{derived-lemma-morphisms-lift-projective}.
We will show below that it follows from (1), (2), (3)
that $K^i \to A^i$ is the inclusion of a direct summand
for every $i$. By property (2) we see that $P^i = \Coker(K^i \to A^i)$
is in $\mathcal{P}$. Thus we can take
$P^\bullet = \Coker(K^\bullet \to A^\bullet)[-1]$ to conclude.

\medskip\noindent
To finish the proof we have to show the following: Let $f : P_1 \to P_2$,
$P_1, P_2 \in \mathcal{P}$ and $P_1/IP_1 \to P_2/IP_2$ is split
injective with cokernel of the form $P_3/IP_3$ for some
$P_3 \in \mathcal{P}$, then $f$ is split injective.
Write $E_i = P_i/IP_i$. Then $E_2 = E_1 \oplus E_3$.
Since $P_2$ is projective we can choose a map $g : P_2 \to P_3$
lifting the map $E_2 \to E_3$. By condition (3) the map $g$
is surjective, hence split as $P_3$ is projective. Set $P_1' = \Ker(g)$
and choose a splitting $P_2 = P'_1 \oplus P_3$. Then $P'_1 \in \mathcal{P}$
by (2). We do not know that
$g \circ f = 0$, but we can consider the map
$$
P_1 \xrightarrow{f} P_2 \xrightarrow{projection} P'_1
$$
The composition modulo $I$ is an isomorphism. Since $P'_1$ is
projective we can split $P_1 = T \oplus P'_1$. If $T = 0$, then
we are done, because then $P_2 \to P'_1$ is a splitting of $f$.
We see that $T \in \mathcal{P}$ by (2).
Calculating modulo $I$ we see that $T/IT = 0$.
Since $0 \in \mathcal{P}$ (as the summand of any $P$ in $\mathcal{P}$)
we see the map $0 \to T$ is surjective and we conclude that $T = 0$
as desired.
\end{proof}

\end{document}
